\section{Completion, finite coefficients, and the arithmetic action}
\label{sec:E}
\label{sec:E-manuscript}

We work in characteristic zero.  In this section $\overline{\mathcal M}_{g,n}$
denotes the stable compactification, while
$\mathcal M_{g,n,\overline{\mathbb Q}}$ denotes geometric base change of
the smooth moduli stack.  This distinction is fixed by the maps in
\S\ref{sec:E-manuscript}: the bar in cote 143 is geometric base change.
Write $\Gamma_{g,n}$ for the orientation-preserving mapping class group
fixing the $n$ labels, and $C_{g,n}$ for a finite-object skeleton of its
boundary exit category.  Arrows point from a nodal stratum towards a
smoothing.  All actions and homomorphisms between profinite groups below
are continuous.

\subsection{Finite congruences and the comparison with sheaves}

\begin{result}{E.1}{Proved}{The strict completion and its finite representations}
For every stable $(g,n)$, set
\[
 \widehat C_{g,n}=\varprojlim_R C_{g,n}/R,
\]
where $R$ ranges over finite-index congruences on all Hom-sets, with the
object set fixed.  Composition makes this a category internal to
profinite sets, and restriction along $C_{g,n}\longrightarrow\widehat
C_{g,n}$ is an equivalence
\[
 \operatorname{Fun}_{\mathrm{cts}}(\widehat C_{g,n},\mathrm{FinSet})
       \simeq \operatorname{Fun}(C_{g,n},\mathrm{FinSet}).       \tag{E.1}
\]
Every functor between two such finite-object categories extends to their
strict completions.  A natural transformation extends as well, and the
extensions preserve its equations.
\end{result}
\begin{proof}
There are finitely many stable labelled dual graphs of a fixed type:
the number of edges is at most $3g-3+n$, the number of vertices is at
most $2g-2+n$, and every vertex genus and every incidence set is bounded.
Jansen's Construction 3.18 identifies isomorphism classes of exit objects
with these graphs.  Choose one object in each class.

For a finite congruence $R$, the quotient keeps identities and all
composition laws.  Intersection of two finite congruences is finite,
because its classes inject into the product of their finite class sets.
The inverse system is therefore cofiltered.  Taking inverse limits of
the finite composition functions defines continuous composition, and
associativity follows in every finite coordinate.

For $F:C_{g,n}\to\mathrm{FinSet}$, equality of the functions $F(a)$ and
$F(b)$ defines a finite congruence on each Hom-set.  There are finitely
many objects, so only finitely many finite function sets occur.  Thus
$F$ factors through one quotient.  Conversely, a continuous map from a
profinite Hom-set to a finite function set factors through finitely many
coordinates; a common refinement handles all pairs of objects.  All
transition maps and projections are surjective on Hom-sets, so the
identity and composition equations descend to this common finite stage.
The components of a natural transformation are finitely many functions;
its naturality equations can be checked on the dense original Hom-sets
or at a common finite stage.  This proves full faithfulness and essential
surjectivity in (E.1).

For a functor $u:C\to D$, pull back every finite congruence on $D$ along
$u$.  The resulting finite congruences on $C$ define the extension
$\widehat u$.  This construction respects composition of functors.
Its application to the finitely many components of a natural
transformation gives the last assertion.  No identification of
$\widehat\Gamma_{g,n}$ with its procongruence image occurs in this proof.
\end{proof}

\begin{result}{E.2}{Known}{Constructible finite coefficients in every stable type}
Fix $\overline{\mathbb Q}\hookrightarrow\mathbb C$.  For every stable
$(g,n)$ there are equivalences, compatible with stalk functors,
\[
 \operatorname{Cons}_{\partial,\mathrm{fin}}
       ((\overline{\mathcal M}_{g,n})_{\overline{\mathbb Q},\mathrm{et}})
 \simeq
 \operatorname{Cons}_{\partial,\mathrm{fin}}
       ((\overline{\mathcal M}_{g,n})^{\mathrm{an}})
 \simeq
 \operatorname{Fun}_{\mathrm{cts}}(\widehat C_{g,n},\mathrm{FinSet}).
                                                               \tag{E.2}
\]
Here $\operatorname{Cons}_{\partial,\mathrm{fin}}$ means finite-stalk
sheaves locally constant on every topological-type stratum of the fixed
boundary stratification.  Jansen, Corollary 6.6 and Theorem 6.7, and
Barwick--Glasman--Haine, Proposition 12.6.4, give the inputs.
\end{result}
\begin{proof}[Derivation, including the stack descent]
The compactification is a finite-type separated Deligne--Mumford stack.
Choose a finite-type quasi-compact etale atlas $U\to\overline{\mathcal
M}_{g,n}$.  Its iterated fibre products form a simplicial finite-type
scheme after refinement by affine etale covers.  On each such scheme,
Proposition 12.6.4 of \cite{ET_BGH2020} identifies constructible etale
sheaves with constructible analytic sheaves.  The comparison is inverse
image, so its compatibility with every face, degeneracy, and refinement
map is part of the construction.  Descent data and their cocycle
equations are consequently identified.  Taking the descent limit gives
the asserted equivalence for the stack; this is also the construction
of \cite[Proposition 13.9.17; Corollary 13.9.19]{ET_BGH2020}.
Restriction to a stratum commutes with inverse image.  Local constancy
and finiteness are etale-local properties.  Hence the equivalence
restricts to exactly the fixed boundary-constructible subcategories.

For extension from $\overline{\mathbb Q}$ to $\mathbb C$, use invariance
of finite etale covers of finite-type schemes under extension of
algebraically closed characteristic-zero fields.  Apply it on every
stratum and on the normal-crossings punctured local charts; the latter
have finite monodromy described by finite quotients of their meridian
lattices.  The specialization maps are maps of finite sets commuting
with these monodromies.  The same invariance thus identifies the
gluing data, and etale descent gives the stack assertion.

For the analytic term, Jansen's Corollary 6.6 identifies the exit
infinity-category with the ordinary Charney--Lee category, in the
variance of Construction 3.18.  Her Theorem 6.7 supplies exodromy.
Take set-valued objects, then the full subcategory with finite values.
Passing to the chosen skeleton is an equivalence; (E.1) completes the
derivation.  We have used the stable-type case of Corollary 6.6 directly,
so no unstable-type inequality printed in a subsequent specialization
is needed.
\end{proof}

\begin{result}{E.3}{Refuted}{The proposed nonsolvable-centralizer obstruction to (E.2)}
The claim in Part E(a) that the finite constructible-sheaf comparison
requires
$C_{\widehat\Gamma}(H)=\overline{C_\Gamma(H)}$ for the Klein-quartic
isotropy group $H\cong\operatorname{PSL}_2(\mathbb F_7)$ is false as
a claimed necessary step.  Formula (E.2) holds there, with this isotropy
retained, by its descent proof.  It does not assert the displayed
centralizer identity.
\end{result}
\begin{proof}
In an etale local quotient chart $[V/H]$, a finite constructible sheaf
is an $H$-equivariant finite constructible sheaf on $V$: the equivariance
maps are indexed by the actual finite set $H$ and satisfy its
multiplication table.  Analytification preserves those maps and their
equations, and its inverse on constructible sheaves transports them
back.  This proves the equivalence on $[V/H]$ for every finite $H$,
including the indicated group of order $168$.  The descent limit
used in E.2 then glues these comparisons.  Neither functor in this
explicit comparison takes a centralizer in a completed mapping class
group.  The surviving action of $H$ on the stalk over a fixed point
is precisely its stack isotropy action.  This proves the positive
comparison at the point singled out in the prompt and removes the
asserted dependency by an explicit construction.
\end{proof}

\begin{result}{E.4}{Proved}{Functorial arithmetic action and its faithfulness}
Forgetting, clutching, and relabelling, with their specified
compatibility isomorphisms, induce functors on the strict completions.
The full completed system has a coherent continuous outer action of
$G_{\mathbb Q}=\operatorname{Gal}(\overline{\mathbb Q}/\mathbb Q)$,
and the map to equivalence classes of its autoequivalences is injective.
At a rational normal-crossings tangential point, a compatible local
representative acts on the meridian lattice by
\[
 (t_e)_{e\in E(G)}\longmapsto
                 (t_e^{\chi(\sigma)})_{e\in E(G)} .           \tag{E.3}
\]
Changing geometric paths conjugates the simultaneous local action.
\end{result}
\begin{proof}
Apply E.1 to the discrete geometric functors and their coherence maps
constructed in Parts A and C.  Alternatively, use (E.2): the algebraic
maps are defined over $\mathbb Q$, and inverse image preserves finite
boundary-constructible sheaves because each source stratum maps into
one target stratum.

For $\sigma\in G_{\mathbb Q}$, the functor $(\sigma^{-1})^*$ on
geometric sheaves gives a left action.  Composition of inverse image
supplies its coherent associativity.  Every finite cover, finite stalk,
and finite gluing diagram used by a constructible sheaf descends to a
number field, so this action is continuous on finite data.  Transferring
through (E.2), including its stalk functors, gives the outer categorical
action.  Compatibility with the geometric operations follows from
their being defined over $\mathbb Q$.

For a smoothing parameter $q_e$, the covers $q_e^{1/m}$ have inertia
$\mu_m$.  The action on these roots of unity is $\chi\bmod m$.
Taking all $m$ proves (E.3), with a single exponent on the whole
local lattice.  A changed path conjugates the inertia subgroup.

For faithfulness, restrict the system to the open stratum of
$\overline{\mathcal M}_{0,4}$, namely
$\mathbb P^1-\{0,1,\infty\}$, with group $\widehat F_2$.
An outer-trivial arithmetic action fixes every conjugacy class of open
subgroups, hence every connected finite cover of the tripod and its
smooth compactification.  Belyi's theorem gives such a presentation
for every elliptic curve over $\overline{\mathbb Q}$.  The action
therefore fixes every elliptic $j$-invariant.  Every algebraic number
is an elliptic $j$-invariant, so the Galois element is the identity.
This argument concerns the full system.  Its $(0,3)$ labelled component
alone is terminal and has the trivial arithmetic action.
\end{proof}

\subsection{The five-point higher comparison}

\begin{result}{E.5}{Proved}{Primitive meridians and retracts at all five-point flags}
Let $G=\pi_1(\mathcal M_{0,5})$.  For a multicurve $\sigma$ on the
labelled five-punctured sphere put $H_\sigma=\operatorname{Stab}_G(\sigma)$.
Every $H_\sigma$ is a retract of $G$.  For $|\sigma|=1$ it is
$\mathbb Z\times F_2$; for $|\sigma|=2$ it is $\mathbb Z^2$.
Every full flag of boundary strata has a connected lifting groupoid,
whose automorphism group is one of
\[
 G,\quad F_2,\quad \mathbb Z\times F_2,\quad
 \mathbb Z^2,\quad\mathbb Z,\quad 1.                         \tag{E.4}
\]
All these groups are good, and the finite categorical quotients induce
their full profinite topologies.
\end{result}
\begin{proof}
Normalize labels $1,2,3$ to $0,1,\infty$, and write the remaining
coordinates as $x,y$.  The open moduli space is the complement of
the five lines
$x=0,x=1,y=0,y=1,x=y$ in $\mathbb C^2$.
Their positive meridians form a basis of $H_1(G,\mathbb Z)=\mathbb Z^5$:
small transverse discs give the generators, and integration of the
five logarithmic differentials gives the dual integral homomorphisms.
The ten boundary meridians have the following columns in this basis:
\[
\begin{array}{c|rrrrr}
12&-1&-1&-1&-1&-1\\
13&0&1&0&1&1\\
14&1&0&0&0&0\\
15&0&0&1&0&0\\
23&1&0&1&0&1\\
24&0&1&0&0&0\\
25&0&0&0&1&0\\
34&-1&-1&0&0&-1\\
35&0&0&-1&-1&-1\\
45&0&0&0&0&1
\end{array}                                                   \tag{E.5}
\]
These are valuations of $x,x-1,y,y-1,x-y$ along the boundary.
For example, at $D_{12}$ put $x=q^{-1}$ and $y=a q^{-1}$ with
$a\notin\{0,1\}$; all five valuations are $-1$.
At $D_{13}$ take $x=1+q,y=1+aq$, and at $D_{23}$ take
$x=q,y=aq$.  The two remaining infinity divisors give the rows
$34,35$ directly.

For a single curve $c$ enclosing a labelled pair $\{a,b\}$,
the cutting sequence is
\[
1\longrightarrow\langle t_c\rangle\longrightarrow H_c
 \xrightarrow{r}\Gamma_{0,4}\longrightarrow1.               \tag{E.6}
\]
Here $r$ is the restriction of the actual forgetting homomorphism
$G\to\Gamma_{0,4}$ obtained by forgetting $a$; on the normalized
boundary it replaces the node by the remaining mark $b$.
Each row in (E.5) is primitive, so there is an integral character
$\ell:G\to\mathbb Z$ with $\ell(t_c)=1$.
The restriction of $(\ell,r):G\to\mathbb Z\times F_2$ to $H_c$
is injective: its kernel lies in $\langle t_c\rangle$ by (E.6),
and $\ell$ kills no nonidentity power.  It is surjective: lift the
$F_2$ coordinate using (E.6) and correct the integer coordinate by
a power of $t_c$.  Its inverse on $H_c$, composed with $(\ell,r)$,
is the required retraction $G\to H_c$.

A two-curve stratum has all three normalized components of type
$(0,3)$, so its stabilizer is the free lattice generated by its two
twists.  Its labels are disjoint pairs, giving exactly fifteen
possibilities.  Every corresponding pair of rows of (E.5) has a
$2\times2$ minor of determinant $1$ or $-1$.  The program
\texttt{code/E\_meridian\_pairs.py} checks all fifteen and writes an
explicit integral left inverse for each two-column matrix.  Composing
this left inverse with $G\to H_1(G,\mathbb Z)$ gives a retraction
onto that twist lattice.  Conjugating the chosen multicurve conjugates
the retraction, so this proves the assertion for every representative.

Consider now a strict flag represented, after individual choices of
transporters, by
$\sigma_0\supsetneq\cdots\supsetneq\sigma_r$.
Pure labelled sphere mapping classes fix each edge of a stabilized
graph, so $H_{\sigma_0}\subset H_{\sigma_i}$ for every $i$.
Successive changes of target trivialization bring every lifted flag to
the canonical one.  Naturality of an automorphism along its arrows
then gives exactly
\[
             \operatorname{Aut}(\text{flag})
                =H_{\sigma_0}/\Delta_{\sigma_r}.             \tag{E.7}
\]
There is no assumption that all representatives for all flags can be
chosen simultaneously as inclusions.  Formula (E.7) is invariant under
the transporters used for one flag.  The possible cardinalities are
$0,1,2$, and (E.6) gives exactly (E.4).

A retract inherits the full profinite topology: any finite quotient of
the retract extends along its retraction.  The indicated lattice
quotients also inherit their quotient profinite topology.  For goodness,
the projection $(x,y)\mapsto x$ is the configuration fibration with
base $\mathbb C-\{0,1\}$ and fibre
$\mathbb C-\{0,1,x\}$, so
$G=F_3\rtimes F_2$ (the group extension splits because its quotient
is free).  Finitely generated free groups and free abelian groups are
good.  For completeness, their free or Koszul resolutions compare
ordinary and continuous finite-coefficient cohomology.  In the free
extension, characteristic finite-index subgroups of $F_3$ are cofinal;
the action on each finite quotient has finite image.  This proves the
exactness of the completed split extension and continuity of the
finite cohomology modules.  Comparing its two Lyndon--Hochschild--Serre
spectral sequences proves goodness of $G$.  The same argument proves
goodness for the direct products in (E.4).
\end{proof}

\begin{result}{E.6}{Proved}{LM3b H.9 at five points}
The nerve of the strict finite-congruence completion of $C_{0,5}$
represents the profinite stratified homotopy type of
$\overline{\mathcal M}_{0,5}$ with its boundary stratification.
In particular it classifies the corresponding constructible sheaves
with $\pi$-finite values, with all higher gluing retained.
\end{result}
\begin{proof}
We verify all full flags, including point--divisor--open flags.
Every component of a genus-zero multicurve determines a split of
the puncture labels.  Distinct disjoint components have distinct
splits: two disjoint curves giving the same split cobound a
puncture-free annulus and are isotopic.  A pure mapping class
preserves each split.  Thus, for fixed representatives $\sigma,\tau$,
every transporter $g$ with $g\sigma\supseteq\tau$ has the same
subset $g^{-1}\tau\subseteq\sigma$.  Consequently
\[
 \{g\in G:g\sigma\supseteq\tau\}=H_\tau g_0
\]
for any one transporter $g_0$.  This is transitivity on the actual
incidence set, not merely on abstract graph types.  It also proves
$H_{\sigma_0}\subseteq H_{\sigma_i}$ for an actual nested flag:
a pure stabilizer fixes each component individually.  The same
argument on the cut components shows that each flag of fixed
labelled graph types has one $G$-orbit.  The labelled splits specify
the retained edges at every step, and homeomorphisms of the
complementary punctured spheres identify any two realizations.

For $N\lhd G$ of finite index define the finite quotient category
using the images in $G/N$ of the transporter sets and of the source
twist subgroups.  These categories are cofinal among the finite
congruence quotients.  Indeed, a functor to a finite category has a
finite-index kernel on each target automorphism group.  Pull it back
to $H_\tau$.  By E.5 there is $N_\tau\lhd G$ of finite index whose
intersection with $H_\tau$ lies in this kernel.  Intersect over the
finitely many targets.  Two transporters with equal image modulo
$N\Delta_\sigma$ differ, after a source twist, by an element of
$N\cap H_\tau$: the preceding fixed-incidence calculation shows
that their ratio fixes $\tau$.  They therefore have equal functor
images.  This
constructs the desired factorization.

For a flag as in (E.7), the lifting groupoid at this finite stage is
transitive, with automorphism group
\[
 p_N(H_{\sigma_0})/p_N(\Delta_{\sigma_r}).                    \tag{E.8}
\]
This follows by the same successive target changes of trivialization.
Every finite-stage arrow lifts to an original transporter arrow
between the chosen objects, so a finite-stage chain has simultaneous
composable lifts.  The transitivity just proved passes to this
lifting groupoid.  We do not assume that an arbitrary quotient of
a curve complex retains all its incidence information.
The nested inclusions $H_{\sigma_0}\subset H_{\sigma_i}$ ensure that
the final automorphism lies in $p_N(H_{\sigma_0})$; its ambiguity is
exactly $p_N(\Delta_{\sigma_r})$.  Thus no extra intersection or
fibre-product assertion is being inferred from a single crossing.

By E.5, system (E.8) is cofinal in all finite quotients of (E.7),
and its group is good.  Goodness for every finite local system,
together with the identical profinite fundamental group, proves that
the corresponding system of classifying spaces is the profinite
homotopy completion of the original flag groupoid.  One may verify
this last implication by obstruction theory into a finite Postnikov
tower: each successive obstruction and choice is a finite local-system
cohomology group, for which goodness is precisely the required
comparison.

All these maps are induced by restriction of actual chains, so they
commute with every face map.  The spatial decollage theorem
\cite[Theorem 2.7.4; Construction 2.8.8 and (2.8.9)]{ET_BGH2020}
reconstructs the profinite stratified object from this full flag
diagram.  The comparison is an equivalence at every entry, hence is
an equivalence of the reconstructed objects.  Finally Corollary
12.6.7 of that source identifies the latter object with the
stratified etale type.  We have checked the three-stratum flags
explicitly; the rank-one argument of LM3b H.8 alone would not give
this step.
\end{proof}

\subsection{The arithmetic subgroup selected by a real framing}

\begin{result}{E.7}{Known}{The AA selection in ANABELIAN E.14}
Fix the standard embedding $G_{\mathbb Q}\hookrightarrow\widehat{
\operatorname{GT}}$ and its complex conjugation $c=(-1,1)$.
Among the maximal closed subgroups satisfying the AA-property of
Hoshi--Mochizuki--Tsujimura, Definition 5.12, exactly one contains $c$;
it is the specified arithmetic image.  This statement retains all
prerequisites of that definition, including its configuration-space
compatibility.  It confirms the subgroup-selection assertion of
ANABELIAN E.14.
\end{result}
\begin{proof}[Derivation from the opened sources]
Theorem 5.17(ii) of \cite{ET_HMT2025} identifies the maximal AA
subgroups with the conjugates of the arithmetic image $G$; its
proof uses Theorem 5.15(iv), Remark 5.12.1, and Lemma 5.16.
Proposition 4(ii) of \cite{ET_LS1997} gives
$C_{\widehat{\operatorname{GT}}}(c)=\langle c\rangle$.
We checked its proof on printed p.~581: commutation forces the
second GT coordinate to be fixed by simultaneous inversion, its
fixed subgroup is trivial, and the hexagon gives multiplier $\pm1$.

Let $bGb^{-1}$ contain $c$.  Then $b^{-1}cb$ is a nontrivial
involution in $G$.  All such involutions in $G_{\mathbb Q}$ are
conjugate: their fixed fields are real closures, the field $\mathbb Q$
has a unique ordering, and an isomorphism of real closures extends
to their algebraic closures.  Thus $b^{-1}cb=g c g^{-1}$ for some
$g\in G$, whence $bg\in C_{\widehat{\operatorname{GT}}}(c)
\subset G$.  It follows that $b\in G$ and $bGb^{-1}=G$.
Existence follows because the arithmetic image itself satisfies AA.

This is a test on closed subgroups of the completed genus-zero
configuration system; Hoshi--Mochizuki--Tsujimura construct it from
the abstract group of the second tripod configuration space.  It
does not identify every individual element satisfying inertia and
permutation conditions with a Galois element.  Once the subgroup
has been selected, membership in it is the additional test.
\end{proof}

\subsection{A mixed-braid calculation in genus one}

\begin{result}{E.8}{Proved}{The kernel for an automorphism of an extension}
Let $1\to K\to E\to Q\to1$ be an exact sequence of finitely generated
profinite groups, with $Z(K)=Z(Q)=1$, and let
$\rho:Q\to\operatorname{Out}(K)$ be its outer monodromy.  Restriction
and passage to the quotient give an isomorphism
\[
 \ker\bigl(\operatorname{Out}(E;K)\longrightarrow
                    \operatorname{Out}(Q)\bigr)
       \simeq C_{\operatorname{Out}(K)}(\rho(Q)).             \tag{E.8}
\]
The same formula holds after imposing any invariant peripheral
condition on the restriction to $K$, by intersecting the right side
with the corresponding subgroup of $\operatorname{Out}(K)$.
\end{result}
\begin{proof}
There is a canonical isomorphism of profinite groups
\[
 E\longrightarrow
 \{(q,a)\in Q\times\operatorname{Aut}(K):[a]=\rho(q)\},
 \qquad e\longmapsto (\bar e,\operatorname{conj}_e|_K).
                                                               \tag{E.9}
\]
The kernel is $Z(K)$.  To prove surjectivity, lift $q$ to $e\in E$;
then $a(\operatorname{conj}_e|_K)^{-1}$ is inner, so multiplying $e$
by an appropriate element of $K$ gives the prescribed pair.  The
continuous bijection is an isomorphism, since the groups are compact
and Hausdorff.  This proof does not require injectivity of $\rho$.

An automorphism of $E$ inducing the identity on $Q$ restricts to an
outer automorphism centralizing $\rho(Q)$.  Conversely, for a lift
$a\in\operatorname{Aut}(K)$ of such an outer automorphism, the map
$(q,b)\mapsto(q,aba^{-1})$ of the right side of (E.9) defines an
automorphism of $E$ over $Q$.  Changing $a$ by an inner automorphism
changes this lift by an element of $\operatorname{Inn}(K)$.
If an automorphism over $Q$ is inner in $E$, its conjugating element
maps to $Z(Q)=1$, and hence belongs to $K$.  Finally, an element of
the kernel on the left of (E.8) can be modified by an inner
automorphism of $E$ to induce the identity on $Q$.  These operations
are inverse and prove the assertion, including its peripheral version.
\end{proof}

\begin{result}{E.9}{Proved}{An explicit inertia-preserving GT lift to
the three-point elliptic group}
Put $K=\widehat\pi_{1,2}$, $Q=\widehat\Gamma_{1,2}$ and
$E=\widehat\Gamma_{1,3}$, where all marks are ordered and the
map $E\to Q$ forgets the third mark.  The standard GT action on
$\widehat B_5$ induces a continuous homomorphism
\[
 \Psi:\widehat{\operatorname{GT}}\longrightarrow
           \operatorname{Out}(E;K)
                                                               \tag{E.10}
\]
whose quotient action on $Q=\widehat B_4/Z(\widehat B_4)$ is the
standard action of Minamide--Nakamura, Theorem~4.3(i).  Its restriction
to $K$ fixes both cusp classes with the common multiplier $\lambda$,
and its action on $E$ preserves nonseparating twist inertia.
\end{result}
\begin{proof}
We specify the quotient, its topology, and the action.  Let
$B_{4,1}\subset B_5$ be the subgroup whose permutation fixes strand
$5$.  Forgetting strand $5$ splits by the first four strands and gives
\[
 B_{4,1}=F_4\rtimes B_4,
 \quad F_4=\langle x_1,x_2,x_3,x_4\rangle,
 \quad x_i=x_{i5}.
\]
Set $c=x_1x_2x_3x_4$ and $\omega_4=(\sigma_1\sigma_2\sigma_3)^4$.
Our Artin convention is
\[
 \sigma_i(x_i)=x_ix_{i+1}x_i^{-1},\qquad
 \sigma_i(x_{i+1})=x_i,
\]
with the other generators fixed.  Thus every $\sigma_i$ fixes $c$,
and $\omega_4$ acts by $\operatorname{Inn}(c)$.  These identities
follow by successively applying the displayed formulas; the exact
reduced-word checks for four, and also three, generators are recorded
in \texttt{code/E\_braid\_models.py} and its output.

Complete this split sequence.  Its kernel has the full profinite
topology: characteristic finite-index subgroups of $F_4$ are cofinal,
and the action on each corresponding finite quotient factors through
a finite quotient of $B_4$.  Hence
$\widehat B_{4,1}=\widehat F_4\rtimes\widehat B_4$.
Write
\[
 D=\widehat F_4/\overline{\langle\!\langle x_i^2:1\leq i\leq4
                              \rangle\!\rangle}
       =C_2\amalg C_2\amalg C_2\amalg C_2,
 \qquad K_4=\ker(D\to C_2),                               \tag{E.11}
\]
where every factor maps nontrivially.  We use the same letters for
the images of $x_i,c$.  The discrete index-two covering is the double
cover of a disc branched at four points.  Its surface has genus one
and two boundary components.  Replacing those boundaries by punctures
identifies $K_4$ with $K$; the two peripheral classes are generated by
$c$ and $x_1cx_1^{-1}$.  The Schreier basis $x_1x_i$, $2\leq i\leq4$,
also proves directly that $K_4$ is free profinite of rank three.

Lifting a half-twist of two adjacent branch points to this double
cover gives a Dehn twist on a nonseparating curve.  Thus the Artin
action on $K_4$ is the based version of elliptic universal monodromy.
After removing base points its quotient is the familiar geometric
isomorphism
\[
 B_4/\langle\omega_4\rangle\simeq\Gamma_{1,2};
\]
this is precisely the isomorphism used in
\cite[Corollary C and the preceding paragraph]{ET_MN2022}.
The full profinite center of $\widehat B_4$ is
$\overline{\langle\omega_4\rangle}\simeq\widehat{\mathbb Z}$
\cite[\S2, equations (2.2) and (2.9)]{ET_MN2022}.
Consequently the element $(c^{-1},\omega_4)$ is central in
$K_4\rtimes\widehat B_4$ and the quotient
\[
 E'=(K_4\rtimes\widehat B_4)\big/
       \overline{\langle(c^{-1},\omega_4)\rangle}          \tag{E.12}
\]
fits into $1\to K_4\to E'\to Q\to1$.  Its central subgroup has
trivial intersection with $K_4$, as is seen by projection to
$\widehat B_4$.  Its outer monodromy is the universal monodromy just
identified.  The geometric universal curve gives the exact sequence
$1\to K\to E\to Q\to1$; this is also the exact sequence in
\cite[Theorem 6.13]{ET_HMTopicsI2012} with all primes retained.
Applying the explicit isomorphism (E.9) to both extensions gives an
isomorphism $E'\simeq E$ inducing the specified identity on kernel and
quotient.  In particular, we have identified the extension itself,
and have not inferred its class from an abstract isomorphism of
the two endpoint groups.

Let $F=(\lambda,f)\in\widehat{\operatorname{GT}}$.  Its genuine
braid action is
\[
 \sigma_1\mapsto\sigma_1^\lambda,
 \qquad
 \sigma_i\mapsto
 f(y_i,\sigma_i^2)^{-1}\sigma_i^\lambda f(y_i,\sigma_i^2),
 \quad y_i=\sigma_{i-1}\cdots\sigma_1^2\cdots\sigma_{i-1}.
                                                               \tag{E.13}
\]
These are \cite[\S0.1, equation (1)]{ET_HS2000}.  They preserve
the permutation quotient and the subgroup on the first four
strands.  They also preserve $\widehat F_4$: the generators $x_{i5}$
are taken to conjugates of their $\lambda$-powers, with conjugators
in the pure subgroup of $\widehat B_5$.  That subgroup lies in
$\widehat B_{4,1}$, where the normal closure of these generators is
$\widehat F_4$.  Applying the inverse action gives equality.
To check conjugacy inside $\widehat F_4$, write such a conjugator
as $vb$, with $v\in\widehat F_4$ and $b\in\widehat P_4$.
The pure Artin action of $b$ sends $x_i$ to an $\widehat F_4$-conjugate
of $x_i$.  This proves the required internal conjugacy as well.

It follows that the closed normal subgroup generated by $x_i^2$
is preserved.  Since $\lambda$ is odd, the parity homomorphism in
(E.11) is preserved too.  The appendix of \cite{ET_HS2000},
printed pp.~34--35, proves the genuine identities
$F(y_i)=y_i^\lambda$ and $F(\omega_i)=\omega_i^\lambda$.
Here $c=y_5$, so both $c$ and $\omega_4$ are taken to their
$\lambda$-powers.  Their commuting quotient relation in (E.12)
is consequently preserved, with
$(c^{-1},\omega_4)\mapsto(c^{-1},\omega_4)^\lambda$.
We obtain an action on $E'$, hence (E.10), preserving every equation
of the extension.  On its quotient $Q$ the action is exactly (E.13)
for four strands, which proves the asserted comparison map.

The first cusp generator $c$ is taken exactly to $c^\lambda$;
the second is conjugated by an odd element of $D$, and thus remains
in the second $K_4$-conjugacy class with exponent $\lambda$.
Finally, the element $\sigma_1$ in (E.12) is the lift of the
half-twist supported on a small two-branch-point disc.  Its annular
double cover avoids the extra base point, so this lift is a
nonseparating Dehn twist on the three-point torus.  Equation (E.13)
takes it to its $\lambda$-power.  All nonseparating curves are
conjugate under the pure mapping class group of this surface:
cutting along either gives a connected genus-zero surface with the
same three ordered punctures and two ordered boundary components,
and an orientation-preserving identification glues back.  Thus all
their completed inertia subgroups are preserved.  This finishes
the construction without using the high-genus extension argument
of Boggi's Lemma~6.16.
\end{proof}

\begin{result}{E.10}{Proved}{Automatic genus-one inertia: LM3b E.9}
Let $A_{\rm per}$ be the subgroup of
$\operatorname{Out}(\widehat\Gamma_{1,3})$ preserving the kernel
$K=\widehat\pi_{1,2}$ of forgetting the third mark and permuting its
two peripheral subgroup conjugacy classes.  Then
\[
 A_{\rm per}=\langle\delta\rangle\times
                    \Psi(\widehat{\operatorname{GT}})
       \simeq C_2\times\widehat{\operatorname{GT}},          \tag{E.14}
\]
where $\delta$ interchanges the two original labels.  Every element
has one common peripheral exponent, normalizes $\langle\delta\rangle$,
and preserves nonseparating twist inertia.  Thus the group $A$ and
its subgroup $A^I$ in LM3b E.1/E.9 coincide, and
$A/\mathfrak S_2\simeq\widehat{\operatorname{GT}}$.
\end{result}
\begin{proof}
The quotient group $Q=\widehat B_4/Z(\widehat B_4)$ is center-free
\cite[\S2, equation (2.9)]{ET_MN2022}; $K$ is free profinite of
rank three and is center-free.  By (E.8), the kernel of
$A_{\rm per}\to\operatorname{Out}(Q)$ is
\[
 C_{\operatorname{Out}^{C}(K)}(\rho(Q)).                    \tag{E.15}
\]
Here $\operatorname{Out}^{C}$ means permutation of the two
peripheral subgroup classes, without a prescribed exponent.
\cite[Theorem 6.13(ii), equivalently Theorem D(i)]{ET_HMTopicsI2012}
identifies (E.15) with the automorphism group of the tautological
twice-punctured elliptic curve.  Its hypotheses are satisfied:
characteristic zero, all primes, $(g,r)=(1,2)$, $2g-2+r=2>1$,
and the open subgroup there is the whole $Q$.  The theorem and its
proof were checked: for $r>0$ its claims $(\ast1)$ and $(\ast2)$
reduce to Theorem~6.12(i); that theorem reduces to the two-point
configuration case and Lemmas~6.8 and~6.10, and no centralizer
identity for a finite nonsolvable subgroup is used.
The printed induction in Lemma~6.8, p.~132, says $3g-3+r>1$
after the complexity-zero base case.  The argument works for
$3g-3+r>0$.  At the omitted type $(1,1)$ choose the irreducible
one-node stable curve, whose normalization is a tripod.  Inserting
the moving point at its node gives two tripod vertices joined by
two edges.  The two one-edge restrictions in that proof are the
already established tripod case; they kill the two independent
Dehn coefficients, respectively.  At $(0,4)$ use two tripods
joined at a node and the same moving-node insertion.  This supplies
both complexity-one cases, so the application here does not
depend on the erroneous strict inequality.

This tautological automorphism group is $C_2$.  Its nonidentity
element is $z\mapsto p_1+p_2-z$ on the elliptic curve: it exchanges
the two punctures.  In (E.11) its outer action is conjugation by
$x_1$, acting on the Schreier basis $x_1x_i$ by inversion.  Hence
it is nontrivial.  Conjugation by $x_1$ also normalizes
$K_4\rtimes\widehat B_4$ inside $D\rtimes\widehat B_4$ and
preserves the central relation (E.12), so it supplies the same
kernel element in $A_{\rm per}$.  Geometrically it is the
relabeling $\delta$.  Explicitly, conjugation by $x_1$ fixes that
relation because its $D$-component is
\[
 x_1c^{-1}\omega_4(x_1^{-1})
   =x_1c^{-1}(cx_1^{-1}c^{-1})=c^{-1}.
\]
Thus deck descent has been checked on the defining central
generator itself.  As an orientation-preserving
relabeling it
preserves nonseparating twist inertia and has peripheral exponent
one.

By \cite[Theorem 4.3(i), Corollary C]{ET_MN2022},
$\operatorname{Out}(Q)=\widehat{\operatorname{GT}}$ via the
four-strand action.  E.9 gives a section of the map from
$A_{\rm per}$ onto this group.  Its kernel has order two and is
therefore central: conjugation acts through
$\operatorname{Aut}(C_2)=1$.  This proves (E.14).
The section and its kernel both preserve the stated inertia,
and their peripheral exponents are respectively $\lambda$ and
$1$.  The final assertions follow for every element of the direct
product, including the formerly imposed normalizer condition.
\end{proof}

\subsection{Explicit elliptic boundary comparisons}

\begin{result}{E.11}{Proved}{The two elliptic boundary maps and forgetting}
Put $G_2=\widehat B_4/\langle\omega _4\rangle$ and
$G_1=\widehat B_3/\langle\omega _3^2\rangle$, where
$\omega _r=(\sigma _1\cdots\sigma _{r-1})^r$ and brackets on
central profinite subgroups mean closed subgroups.  Under the
geometric identifications $G_i=\widehat\Gamma_{1,i}$, forgetting
the second mark is
\[
 q:G_2\longrightarrow G_1,\qquad
 (\sigma _1,\sigma _2,\sigma _3)\longmapsto(a,b,a),
 \quad a=\sigma _1,\ b=\sigma _2\text{ in }G_1.       \tag{E.16}
\]
The rational-tail boundary has stabilizer $\widehat B_3\subset G_2$,
twist $\omega _3^2$, and normalization group $G_1$.
For the nonseparating twist $a=\sigma _1\in G_2$, put
$v=\sigma _2\sigma _1^2\sigma _2$ and $d=\sigma _3$.
Its stabilizer $H_a$ is generated by $a,v,d$ and satisfies
\[
 [a,v]=[a,d]=1,\qquad (vd)^2=(dv)^2,\qquad
                    a^2(vd)^2=1.                         \tag{E.17}
\]
Consequently
\[
 H_a/\langle a\rangle\simeq
 \widehat A_3/\langle\omega _3\rangle,
 \quad A_3=\langle\sigma _1^2,\sigma _2\rangle\subset B_3,
                                                               \tag{E.18}
\]
with $\sigma _1^2\mapsto v$ and $\sigma _2\mapsto d$.
This is the fundamental group of
$\mathcal M_{0,\{1,2\},[+,-]}$, the four-point normalization
with its two node branches unordered.  Its geometric branch
character is $v\mapsto1$, $d\mapsto0$ in $C_2$.
For every
$F=(\lambda,f)\in\widehat{\operatorname{GT}}$, the genuine
braid actions commute with (E.16), with the rational-tail
quotient, and with (E.18), including its branch permutation.
\end{result}
\begin{proof}
The geometric braid identification and the Cardano--Ferrari
map are checked in \cite[Corollary C; proof of Proposition
3.2(i), especially (3.1)--(3.5)]{ET_MN2022}.  Keeping the central
involution instead of quotienting it gives precisely (E.16).
The one-boundary torus mapping class group is $B_3$; the
two-chain relation identifies its boundary twist with
$(ab)^6=\omega _3^2$.  Gluing its boundary to the tripod
carrying marks $1,2$ gives the rational-tail stabilizer stated
above.  Its map to $G_2$ is injective: cutting the rational-tail
curve has kernel its infinite twist subgroup and quotient
$\Gamma_{1,1}$, and this is the same central extension as
$B_3\to\mathrm{SL}_2(\mathbb Z)$.

For completeness the centralizer presentation comes directly
from cutting the disk braid along a small disk containing its
first two strands.  A braid preserving that disk is determined
by its internal half twist $a$ and a braid on the collapsed disk
and the other two strands.  The latter braid fixes the collapsed
strand and hence belongs to $A_3$.  The framing of that strand
gives the central $a$-factor.  The two generators of $A_3$ lift
to $v,d$; its presentation is
$\langle v,d\mid vdvd=dvdv\rangle$.
Restoring the four-strand full twist gives
$\omega _4=a^2(vd)^2$ and gives (E.17) after quotienting.
These identities are also verified in the faithful Artin action
on the free group of rank four by the attached program
\texttt{E\_braid\_models.py}.
The cut genus-one surface is a sphere with its original two
punctures and two branches.  The disjoint twist $d$ fixes both
branches.  The order-two element $vd$ modulo $a$ exchanges
them, as the half-turn of the cut cylinder.  Thus the branch
character has values $v\mapsto1,d\mapsto0$.  Its kernel is
the pure four-point group.  In the original mixed braid group
$A_3$ the ordinary strand-permutation character instead has
values $v\mapsto0,d\mapsto1$; these two explicit characters
must be retained in using the quotient identification.

These descriptions pass to the full profinite groups.
For genus one the congruence topology is the full topology;
the cutting sequences and the centralizer of each twist are
exactly \cite[Corollary 4.9, Theorem 4.10, Corollary 4.12]{ET_BF2023}.
The factors after cutting have genus zero or one.  The exact
cutting sequences, with the injected completed twist lattice,
identify their completed stabilizers with their closures.
Thus taking a full profinite centralizer adds no elements to
the displayed completed geometric stabilizer.

The genuine action of $F$ on $B_n$ is
\cite[formula (1)]{ET_HS2000}.  It restricts to the standard
action on $B_3\subset B_4$ and sends $\omega _3$ to
$\omega _3^\lambda$.  This proves the rational-tail assertion.
For (E.16), its only additional generator is $\sigma _3$.
Here $y_3=\sigma _2\sigma _1^2\sigma _2=v$ and
\[
 q(v)=ba^2b=-a^{-2}\quad\hbox{in }\mathrm{SL}_2(\mathbb Z).
\]
The equality follows either from the braid presentation and
$(ab)^6=1$, or from the usual faithful integral matrices
$a=\left(\begin{smallmatrix}1&1\\0&1\end{smallmatrix}\right)$,
$b=\left(\begin{smallmatrix}1&0\\-1&1\end{smallmatrix}\right)$.
Thus $q(v)$ commutes with $a^2$.  Since $f$ lies in the
closed commutator subgroup of the free profinite group,
$f(q(v),a^2)=1$.  Formula (1) therefore sends $q(\sigma _3)$
to $a^\lambda$, exactly as required.  It already commutes on
$\sigma _1,\sigma _2$, proving the equality of genuine maps.

Finally \cite[\S0.3, pp.~7--8]{ET_HS2000} states and derives
the genuine braid-tower action which commutes with the
string-doubling homomorphism
\[
 D:\widehat A_3\longrightarrow\widehat B_4,
 \qquad D(\sigma _1^2)=v,\quad D(\sigma _2)=d.
\]
It also preserves the strand permutation.  Since the actions
send $a$ and $\omega _3$ to their $\lambda$th powers, they
descend through (E.17)--(E.18).  Moreover $v$ maps to $v^\lambda$
and $d$ to a conjugate of $d^\lambda$; $\lambda$ is odd.
They therefore preserve the geometric branch character just
computed, as well as the ordinary strand character.
This proves all the asserted comparisons as equalities before
passing to outer automorphisms.
\end{proof}

\begin{result}{E.12}{Proved}{The orbifold normalizer maps and their parameter}
For $r=3,4$, put $F_r=\widehat F(x_1,\ldots,x_r)$ with its
$r+1$ peripheral classes, $D_r=F_r/\overline{\langle\!\langle
x_1^2,\ldots,x_r^2\rangle\!\rangle}$, and
$K_r=\ker(D_r\to C_2)$, every $x_i$ mapping to the
nonidentity element.  Then $K_3=\widehat\pi_{1,1}$ and
$K_4=\widehat\pi_{1,2}$.  Let $\mathcal A_r$ be the genuine
peripheral automorphism group fixing the infinity class, and
$\widetilde N_r$ the inverse image in $\mathcal A_r$ of the
normalizer of the full geometric sphere group
$\Sigma_{0,r+1}\subset\operatorname{Out}_{\rm lac}(F_r)$.
Quotient-and-restriction
\[
 \rho_r:\mathcal A_r\longrightarrow\operatorname{Aut}_{\rm lac}(K_r)
\]
sends $\widetilde N_r$ into the outer normalizer of the full
elliptic geometric group $\Sigma_{1,r-2}$.  It descends to the
normalizer quotients by these geometric groups and is the
identity on their standard $\widehat{\operatorname{GT}}$
parameters.  The analogous elliptic filling map has the same
parameter.  This settles the orbifold-normalizer portion of
LM3b E.4, with its specified maps.
\end{result}
\begin{proof}
A peripheral automorphism takes $x_i^2$ to a conjugate of
$x_{\sigma(i)}^{2\lambda}$ for a unit $\lambda$.
Applying the same assertion to its inverse proves equality
of the closed normal subgroups generated by the squares.
It preserves parity because $\lambda$ is odd.  With
$c=x_1\cdots x_r$ in $D_r$, the cusp of $K_3$ is $c^2$,
and the two cusps of $K_4$ are $c,x_1cx_1^{-1}$.
An odd conjugator is made even by multiplication by $c$ when
$r=3$, and exchanges the two cusp classes when $r=4$.
This proves the asserted genuine peripheral homomorphism.
Changing its genuine representative by $\operatorname{Inn}(g)$
changes its outer restriction by either $1$ or the geometric
deck involution $\delta_r=[\operatorname{conj}_{x_1}]$, according
to the parity of $g$.

Here is the full normalizer check.  The genus-zero universal
curve gives $\widehat\Gamma_{0,r+2}$ as the inverse image of
the pure geometric subgroup of $\operatorname{Out}(F_r)$ in
$\operatorname{Aut}(F_r)$.  Conjugation therefore embeds its
normalizer quotient into
$\operatorname{Out}(\widehat\Gamma_{0,r+2})$.
Indeed a representative centralizing the inner copy of $F_r$
is the identity, since $F_r$ is center-free.
By \cite[Corollaries 2.6(i), 2.8(i)]{ET_HMM2022}, its target is
$\widehat{\operatorname{GT}}\times\mathfrak S_{r+2}$ with no
extra inertia assumption: here the curve is the tripod and its
configuration length is $r-1\geq2$.  The image normalizes the
geometric point-stabilizer $\mathfrak S_{r+1}$.  Its normalizer
inside $\mathfrak S_{r+2}$ is itself, since its unique fixed
letter is preserved.  Dividing by that geometric subgroup gives
an injection into $\widehat{\operatorname{GT}}$.
Conversely the genuine mixed braid action provides every
parameter.  Thus each element of $\widetilde N_r$ is,
modulo inner automorphisms, a geometric infinity stabilizer
times the restriction $\alpha_F$ of the genuine
$(r+1)$-strand action for a unique $F$.
Both factors preserve infinity, so the geometric factor can
indeed be taken in its stabilizer.

The infinity stabilizer is the image of $B_r$.  Lifting the
branched double cover of its disk identifies its action on
$K_r$ with the elliptic geometric group: for $r=3$ it is
$B_3/\langle\omega _3^2\rangle$, and for $r=4$ it is
$B_4/\langle\omega _4\rangle$ together with $\delta_4$.
The first assertion is the torus two-chain relation; the second
is the braid realization in E.9.  The genuine mixed braid
action preserves $F_r\rtimes\widehat B_r$, intertwining the
action of every $b\in\widehat B_r$ on $K_r$ with that of its
standard $F$-image.  It also normalizes the deck group.
These are all the geometric generators, which proves the full
normalizer inclusion.  It proves descent through changes of
inner representative, infinity stabilizer, and sheet as well.

For $r=4$ the induced action on
$B_4/\langle\omega _4\rangle$ has parameter $F$, precisely
the parameter detected in E.9--E.10.  For $r=3$ the same
mixed-braid construction gives
\[
 E_3=(K_3\rtimes\widehat B_3)/
       \langle(c^{-2},\omega _3^2)\rangle
       \simeq\widehat\Gamma_{1,2}.                       \tag{E.19}
\]
To verify the extension identification, the branched disk cover
gives its geometric monodromy on $K_3$; the square of the
full twist acts by $\operatorname{Inn}(c^2)$.  For a
center-free kernel every extension is the pullback of
$\operatorname{Aut}(K_3)\to\operatorname{Out}(K_3)$ along
this monodromy, by the proof of E.8.  This part of that proof
does not require the quotient to be center-free.
It proves (E.19) and the displayed central relation exactly.
The action sends both entries of that relation to their
$\lambda$th powers and hence descends.

The included $B_3$ in (E.19) consists of the two standard
nonseparating twists on the elliptic component separated by
the rational tail.  It is therefore the stabilizer in E.11,
and its action is the genuine standard $F$-action.  This
detects the parameter in
$\operatorname{Out}(\widehat\Gamma_{1,2})=
\widehat{\operatorname{GT}}$.
Indeed an inner conjugacy in $G_2$ between actions preserving
this $B_3$ normalizes its central boundary-twist subgroup and
therefore belongs to $B_3$ by the twist-normalizer theorem.
Restriction thus gives an inner conjugacy in $B_3$ itself.
The standard map
$\widehat{\operatorname{GT}}\to\operatorname{Out}(\widehat B_3)$
is injective: restriction after quotienting its center is the
pair description
$\widehat{\operatorname{GT}}\subset
\widehat{\operatorname{GT}}_0$ of
\cite[\S0.3 and Main Theorem (a)]{ET_HS2000}.
Hence the detected parameter is exactly $F$.
Finally the genuine identity (E.16) proves the same-parameter
assertion for elliptic filling.
\end{proof}

\subsection{The covering-and-filling system}

\begin{result}{E.13}{Proved}{The injectivity asserted in ANABELIAN E.13}
Let $\mathscr H$ be the labelled and coloured category of
ANABELIAN E.8: its objects are oriented hyperbolic surfaces;
its generators are finite unramified covers and specified
puncture fillings; its relations are diagram isotopy and
equality of actual composite maps.  Complete a right mapping-class
orbit $S_u\backslash M_S$ to
$\overline{S_u}\backslash\widehat M_S$ and retain the finite
label-permutation quotients, all label-forgetting maps, and all
specified deletions.  Let $F$ be a continuous type-preserving
autoequivalence preserving this structure.  If its outer action
on every labelled and coloured genus-zero mapping class group
is trivial, its induced outer action on every
$\widehat M_{g,n}$, $2g-2+n>0$, is trivial.

Thus restriction is injective on the compatible families of
outer group automorphisms induced by these autoequivalences.
The statement is about these families, not a classification of
the autoequivalences inducing a given family.  No congruence
subgroup property is used.
\end{result}
\begin{proof}
We give the categorical reconstruction and the required
Hurwitz surjection, since both are essential to the assertion.
There are finitely many right orbits in each uncompleted Hom
set.  A nonidentity elementary arrow strictly decreases
$2g-2+n$, covering degrees are bounded by the ratio of endpoint
complexities, and a finitely generated group has only finitely
many subgroups of bounded index.  Fillings with specified
colours have finitely many orbits.  Composing two finite lists
of orbits still gives finitely many, since the intermediate
mapping class group permutes the first list.  Its stabilizer of
an orbit is consequently open in its full profinite completion.
The completed orbit formula in the statement therefore has
continuous left and right actions and composition; associativity
follows from density.  In particular stabilizers of completed
arrows are exactly closures of the indicated discrete right
stabilizers.  This verifies the feature of E.8 used below.

Put $P_S=\widehat\pi_1(S)$.  Covers give injections of their
specified finite index and fillings give surjections on these
groups, as outer continuous homomorphisms.  These assignments
extend to completed arrows: for finitely generated profinite
groups $P,Q$, continuous homomorphisms $P\to Q$ form a closed
subspace of a finite power of $Q$, and quotienting by target
conjugacy is compact Hausdorff.  Composition is continuous,
as is checked in each finite characteristic quotient.
We use this functor to recognize equations between arrows;
we do not suppose that $F$ already commutes with it.

\emph{Point pushing and chosen sheets.}
For affine hyperbolic $S$, fill one new distinguished puncture
by $w:S_p\to S$.  The topological Birman sequence gives
$1\to\pi_1(S,p)\to M_{S_p}\to M_S\to1$.
Its completed kernel is exactly $P_S$: the composite
\[
 P_S\longrightarrow\widehat M_{S_p}
             \longrightarrow\operatorname{Aut}(P_S)
\]
is the inner-automorphism map and is injective, because a
nonabelian free profinite group is center-free.  Right
exactness of completion then gives
\[
 1\longrightarrow P_S\longrightarrow\widehat M_{S_p}
                \longrightarrow\widehat M_S\longrightarrow1. \tag{E.20}
\]
It is the right stabilizer of $w$, and therefore is preserved
by $F$.  This recovers an outer automorphism $\psi_S$ of $P_S$.
Two-deletion squares show that $\psi_S$ preserves the kernel
$N_c=\overline{\langle\!\langle I_c\rangle\!\rangle}$ of
filling any old labelled puncture $c$ whenever the resulting
surface is still affine hyperbolic.

For a degree-$d$ cover $f:S\to T$, remove an ordinary target
point $q$ and all $d$ points $D=f^{-1}(q)$, distinguishing
$p\in D$.  Write $f':S_D\to T_q$, $u:S_D\to S$ and
$v:T_q\to T$ for the resulting square.  An image of this
square under $F$ is still the correct puncture pullback.
Indeed, for the images $J'\subset P_{T_q}$ and $J\subset P_T$
of its two covering homomorphisms, commutativity up to a
chosen target conjugator gives
\[
 J'\subset v_*^{-1}(J),\qquad
 [P_{T_q}:J']=[P_{T_q}:v_*^{-1}(J)]=d.
\]
The subgroups are equal.  Its $d$ new source cusps therefore
lie over $q$ with degree one; no old target cusp is killed by
$v_*$.  This also identifies which cusps the square deletes.

We need one cancellation calculation.  If $u$ fills extra
ordinary source punctures and $q_u:M_{S_D}\to M_S$ forgets
them, then
\[
 \operatorname{Stab}_{M_{S_D}}(fu)
                  =q_u^{-1}\operatorname{Deck}(f).        \tag{E.21}
\]
For if $fu a=fu$, surjectivity on surface groups implies that
$q_u(a)$ induces the same outer action as an element of
$N_{\pi_1T}(f_*\pi_1S)/f_*\pi_1S$, the deck group.
The injective discrete Dehn--Nielsen--Baer representation
then gives (E.21); the reverse inclusion is immediate.
The finite deck group embeds in $\widehat M_S$.  One direct
verification uses a graph spine of its affine quotient:
\[
 H_1(S,\mathbb Q)\simeq\mathbb Q\oplus
     \mathbb Q[\operatorname{Deck}(f)]^{,r-1},\qquad r\geq2.
\]
The action is faithful and hence is detected modulo some
integer on integral homology.  Since $q_u$ is surjective,
its completed kernel is the closure of its kernel, and the
preimage in (E.21) is a finite union of kernel cosets.  Thus
(E.21) remains exact after completion and in arbitrary
completed markings.

Now characterize inside $P_T\subset\widehat M_{T_q}$ the
elements $b$ admitting a lift $a$ with
\[
                  b f'=f'a,\qquad ua=u.                  \tag{E.22}
\]
The chosen sheet makes the lift unique, because a deck
transformation fixing a sheet is the identity.  Discretely
(E.22) says exactly that the loop $b$ returns that sheet;
its subgroup is $f_*\pi_1(S,p)$.  The same holds after
completion.  The target lifting subgroup is open; on its
intersection with $P_T$, (E.21) says the forgotten lift lies
in the finite deck group.  The second equation in (E.22)
is the kernel of that finite continuous quotient.  Its
discrete subgroup is dense, proving that its closure is
$f_*P_S$.  Filling $D\setminus\{p\}$ identifies its lift
with point pushing in $S_p$.
All of (E.22), its chosen sheet and its fillings are retained
categorical data.  Applying $F$ gives
\[
                  (Ff)_*\psi_S=\psi_T f_*                \tag{E.23}
\]
up to one target inner automorphism.  This proves the needed
naturality rather than imposing it.

\emph{Genus-zero gauges and labelled ramification.}
For a genus-zero surface with $b\geq5$ punctures, every
coloured mapping class completion is center-free.  The pure
case follows from (E.20), starting with the free four-point
group.  For the coloured case its permutation acts faithfully
on the rational abelianization of the pure group:
\[
 H_1(\Gamma_{0,b},\mathbb Q)=
 \mathbb Q^{\binom b2}/
 \left\langle\sum_{j\ne i}e_{ij}:1\leq i\leq b\right\rangle.
                                                               \tag{E.24}
\]
The arrangement presentation with three punctures fixed at
$0,1,\infty$ proves this formula by its independent meridians.
Distinct $e_{ij}$ stay distinct for $b\geq5$: a difference of
two would have coefficients $a_i+a_j$ supported on two edges;
the complement is connected with an odd cycle, forcing all
$a_i=0$.  A permutation acting trivially fixes every pair and
is the identity.  The same argument over $\mathbb Q_\ell$
detects the finite permutation in the profinite group.

The hypothesis on $F$ consequently allows gauges making the
actions on $T,T_q$ and the filling $v$ all identities.
Explicitly, first fix the target action and then $v$ by its
single right orbit.  The remaining inner source conjugator
maps to the center of the target, hence lies in $P_T$;
gauging by its inverse fixes $v$ and the source action.
Thus $\psi_T$ in (E.23) may be taken to be the identity.
It follows that $Ff$ and $f$ have conjugate image subgroups,
so the same unlabelled cover orbit.
For a retained source cusp $c$, preservation of $N_c$ in
(E.23) also gives equality of the target normal closures of
$f_*I_c$ and $(Ff)_*I_c$.  In
$P_T^{\rm ab}=\widehat{\mathbb Z}^{b}/
\widehat{\mathbb Z}(1,\ldots,1)$ these have images
$e\widehat{\mathbb Z}[c_i]$ and
$e'\widehat{\mathbb Z}[c_j]$.
Distinct cusp lines are distinct over every $\mathbb Q_\ell$;
equality forces $i=j$ and then equality of all prime valuations
of the positive integers $e,e'$, so $e=e'$.
In particular a retained simple ramification point over a
distinguished branch value is recovered uniquely.  Passing
through the specified finite labelling cover gives the same
assertion with retained colours.

\emph{A Hurwitz cover surjecting on the higher-genus group.}
Fix $g\geq1$ and a stable marked type $(g,n)$.  Choose
$d\geq2g+2n+5$ and set $b=2d+2g-2$.
Let $\mathcal H$ parametrize degree-$d$ simply branched maps
to $\mathbb P^1$ with $n$ labelled simple ramification points,
their separately labelled values, and the remaining branch
values unordered; divide target coordinates by
$\operatorname{PGL}_2$.
Riemann's covering classification gives a finite etale branch
map $\mathcal H\to\mathcal M_{0,n,[b-n]}$: there are finitely
many transposition monodromy tuples, and distinct moving
branch values uniquely transport each cover.  Connected
transposition monodromy is $S_d$, so the deck group is trivial
for $d\geq3$.

We verify the stronger assertion that
$\mathcal H\to\mathcal M_{g,n}$ is smooth, onto and has
geometrically connected fibers.  For a fixed $(C;P_i)$, let
$D$ vary among reduced degree-$d$ divisors disjoint from the
$P_i$.  This is a smooth irreducible open subset of
$\operatorname{Sym}^d C$.  The spaces $H^0(C,\mathcal O(D))$
form a vector bundle; the $n$ independent conditions
$df(P_i)=0$ cut out a vector subbundle.
The conditions of exact simple poles at $D$, simple critical
points and distinct critical values are nonempty open
conditions.  Indeed Riemann--Roch gives surjection onto all
jets of total length at most $2n+4$ under the bound on $d$.
A multiple unspecified critical point imposes two independent
conditions over a one-dimensional locus.  Two unspecified
critical points with equal values impose three conditions
over a two-dimensional locus.  A coincidence involving one
prescribed critical point imposes two over a one-dimensional
locus; two prescribed values impose one proper hyperplane.
Zero pole coefficients and zero second derivatives at the
$P_i$ each impose proper hyperplanes.  These finitely many
incidence types cannot cover the vector space.
Every resulting function has unique pole divisor $D$.
This is precisely the chart of maps with infinity unramified,
and its fibers are irreducible.  The same bundles and jet
maps work in families, proving smoothness and surjectivity.
Every target admits such a coordinate, whose choices form a
connected open subset of $\operatorname{PGL}_2$; quotienting
coordinates retains geometric connectedness of the fibers.
The argument descends from finite labelling covers for colours.

It follows that $\mathcal H$ is connected and
\[
             \pi_1^{\rm et}(\mathcal H)	woheadrightarrow
                          \widehat M_{g,n}.               \tag{E.25}
\]
Indeed a connected finite etale cover of the base stays
connected: a clopen decomposition of its pullback would,
by connected fibers and the openness of the smooth map,
give a clopen decomposition of the base cover.

Choose a cover in $\mathcal H$, remove all its branch values
and their preimages, and retain the $n$ ramification labels.
This gives $f:S\to T$ and a filling $u:S\to(C\setminus\{P_i\})$.
Here $b\geq5$ and $S$ has $b(d-1)$ punctures, so every
single deletion used above remains affine hyperbolic.
The genus-zero gauges and ramification calculation show
$Ff=f$ after a coloured source gauge.  The open target
lifting subgroup $H_f\subset\widehat M_T$ is
$\pi_1^{\rm et}(\mathcal H)$.  Its lift
$\ell_f:H_f\to\widehat M_S$ is unique, since the deck group
is trivial, and is characterized by $hf=f\ell_f(h)$.
Applying $F$ gives $F_S\ell_f=\ell_f$.
Gauge the filled object so that $F(u)=u$; this uses the single
right orbit of the filling and its surjective forgetting map.
The resulting action on $\widehat M_{g,n}$ fixes the image
of $H_f$, which is all of that group by (E.25).
Its outer class is therefore trivial.  This proves every
stable type in positive genus, including proper curves, and
the genus-zero types are the hypothesis.
\end{proof}

